% Copy-ready correction fragments; NOT a standalone document.
% Prepared 2026-10-07 against ProveIt commit
% c78c7c3dc2fc742a1af9492e47d578b3fb4e01e7.
% Read stieltjes-zeros-distribution-duality-review.md before applying.
% Requires amsmath/amssymb; no private macros or cross-document labels.
% These passages are proposed replacements/insertions, not an automatic patch.

\paragraph{Formal distribution-rank statement.}
For every $q\geq2$ and $k\geq1$, the formal quotient by the specified
distribution relations, with the integral coordinate and lower-order jet
terms treated as background, has residual dimension $\varphi(q)-1$.
An all-parameter proof is given by the character normal form in
\emph{Zeros, Distribution Modules, and Regularized Duality in the Stieltjes
Parameter-Derivative Tower}. This is a rank theorem for specified relations;
it does not assert arithmetic independence of the residual numerical constants
or completeness of all numerical identities.

\paragraph{Numerical evidence for unresolved second-derivative constants.}
The reported $200$-digit integer-relation search found no relation in the
tested basis and search range. These computations motivate treating the six
constants as separate candidate generators. They do not prove linear or
algebraic independence, irrationality, or non-elementarity. A reproducible
record must specify the basis, precision, coefficient-height limit, and
iteration budget.

\paragraph{Explicit regularization at a trivial zero.}
Let $\chi$ be primitive of conductor $q$, let $k\geq1$, and suppose
$\chi(-1)=(-1)^k$. Then
\[
 \frac{L'(k+1,\chi)}{L(k+1,\chi)}
 +\overline{\frac{L''(-k,\chi)}{2L'(-k,\chi)}}
 =\gamma+\log\frac{2\pi}{q}-H_k.
\]
Indeed, dividing $L(-k+t,\overline\chi)$ by its simple zero $t$ gives
\[
 V(t)=L'(-k,\overline\chi)
      +\frac{L''(-k,\overline\chi)}{2}t
      +\frac{L'''(-k,\overline\chi)}{6}t^2+\cdots.
\]
Its first logarithmic derivative therefore contains the factor $1/2$.
At second order the corresponding identity is
\[
 \left(\frac{L''}{L}-\frac{(L')^2}{L^2}\right)(k+1,\chi)
 -\overline{\left\{
 \frac{L'''(-k,\chi)}{3L'(-k,\chi)}
 -\left(\frac{L''(-k,\chi)}{2L'(-k,\chi)}\right)^2\right\}}
 =H_k^{(2)}-\frac{\zeta(2)}2.
\]
For imprimitive characters, the conductor and missing Euler factors must be
included before using the primitive functional equation.

\paragraph{Quarter-argument elimination.}
At $k=2$, the displayed identity eliminates the nonprincipal beta-derivative
coordinate in favor of a negative-order polygamma value. The remaining terms
include principal zeta jets and the explicitly displayed named constants.
The formula $\beta'(-1)=2G/\pi$ reduces to a lower value (Catalan) coordinate;
no elementary-reduction or non-elementarity theorem is asserted for $G$.

\paragraph{Character-coordinate scope.}
For every Dirichlet character modulo $q\geq2$, primitive or not,
\[
 \sum_{a=1}^{q}\chi(a)\zeta(s,a/q)=q^sL(s,\chi).
\]
Differentiation extends the character-coordinate identities to all characters.
Primitivity remains necessary for the unmodified primitive functional equation.

\paragraph{Clausen parity.}
With the convention
\[
 \operatorname{Cl}_{2m}(\theta)=\operatorname{Im}
       \operatorname{Li}_{2m}(e^{i\theta}),\qquad
 \operatorname{Cl}_{2m+1}(\theta)=\operatorname{Re}
       \operatorname{Li}_{2m+1}(e^{i\theta}),
\]
one has
\[
 \operatorname{Cl}_{w}(-\theta)=(-1)^{w+1}
       \operatorname{Cl}_{w}(\theta).
\]
Thus the residual Clausen component of the polygamma Fourier bridge is odd at
even weight and even at odd weight. The weight-three cosine component belongs
to the symmetric tetragamma combination. The weight-two Catalan example is
unchanged.
